\documentclass[17pt,compress]{beamer}
\usepackage{beamerthemesplit}
\mode<presentation>
{
  \usetheme{Warsaw}
  \setbeamercovered{transparent}
  \setbeamertemplate{navigation symbols}{}
}
% Taken from Fernando's slides.
\usepackage{ae,aecompl}
\usepackage[scaled=.95]{helvet}

\usepackage[english]{babel}
\usepackage[latin1]{inputenc}
\usepackage[T1]{fontenc}


\definecolor{blue}{rgb}{0.16,0.32,0.75}
\setbeamercolor{structure}{fg=blue}
\author[FOSSEE]{}
\institute[IIT Bombay]{}
\date[]{}
% \setbeamercovered{transparent}

% theme split
\usepackage{verbatim}
\newenvironment{colorverbatim}[1][]%
{%
\color{blue}
\verbatim
}%
{%
\endverbatim
}%

\usepackage{mathpazo,courier,euler}
\usepackage{listings}
\lstset{language=sh,
    basicstyle=\ttfamily\bfseries,
  showstringspaces=false,
  keywordstyle=\color{black}\bfseries}

% logo
\logo{\includegraphics[height=1.30 cm]{3t-logo.pdf}}
\logo{\includegraphics[height=1.30 cm]{fossee-logo.png}

\hspace{7.5cm}
\includegraphics[scale=0.3]{fossee-logo.png}\\
\hspace{281pt}
\includegraphics[scale=0.80]{3t-logo.pdf}}

\begin{document}

\sffamily \bfseries
\title
[Basic Data-types and Operators]
{Python - Basic Data-types and Operators}
\author
[FOSSEE, IIT - Bombay]
{\small Talk to a Teacher\\{\color{blue}\url{http://spoken-tutorial.org}}\\National Mission on Education
 through ICT\\{\color{blue}\url{http://sakshat.ac.in}} \\[0.5cm]{\tiny Script by: Hardik ghaghada \\ Naration by: Hardik Ghaghada \\ 13 August 2015}}

% slide 1
\begin{frame}
   \titlepage
\end{frame}
%%%%%%%%%%%%%%%%%%%%%%%%%%%%%%%%%%%%%%%%%%%%%%%%%%%%%%%%%%%%%%%%%%%%%%%%%%%%%%%%
\begin{frame}
\frametitle{Objectives}
\label{sec-2}

  In this tutorial we will learn, \pause
\begin{itemize}

\item Datatypes in Python\pause
\label{sec-2_1}%
	\begin{itemize}
	\item Numbers\\
	\label{sec-2_1_1}%
	\item Boolean\\
	\label{sec-2_1_2}%
	\item Sequence\\
	\end{itemize} % ends low level
\end{itemize} % ends low level
\end{frame}
%%%%%%%%%%%%%%%%%%%%%%%%%%%%%%%%%%%%%%%%%%%%%%%%%%%%%%%%%%%%%%%%%%%%%%%%%%%%%%%%
\begin{frame}
\frametitle{Objectives contd..}
\label{sec-2}

\begin{itemize}

\item Operators in Python\pause
\label{sec-2_2}%
\begin{itemize}

\item Arithmetic Operators\\
\label{sec-2_2_1}%
\item Boolean Operators\\\pause
\label{sec-2_2_2}%
\end{itemize} % ends low level

\item Python Sequence Datatypes\pause
\label{sec-2_3}%
\begin{itemize}

\item list
\label{sec-2_3_1}%
\item string
\label{sec-2_3_2}%
\item tuple
\label{sec-2_3_3}%
\end{itemize} % ends low level
\end{itemize} % ends low level
\end{frame}
%%%%%%%%%%%%%%%%%%%%%%%%%%%%%%%%%%%%%%%%%%%%%%%%%%%%%%%%%%%%%%%%%%%%%%%%%%%%%%%%
\begin{frame}
\frametitle{System Specifications}\pause
\begin{itemize}
\item Slackware Linux 14.1\pause
\item \texttt{Python 2.7.5} \pause
\item \texttt{IPython 2.3.0}
\end{itemize}
\end{frame}
%%%%%%%%%%%%%%%%%%%%%%%%%%%%%%%%%%%%%%%%%%%%%%%%%%%%%%%%%%%%%%%%%%%%%%%%%%%%%%%%
\begin{frame}
\frametitle{Numbers}
\label{sec-3}


\begin{itemize}
\item \texttt{int}
\item \texttt{float}
\item \texttt{complex}
\end{itemize}
\end{frame}
%%%%%%%%%%%%%%%%%%%%%%%%%%%%%%%%%%%%%%%%%%%%%%%%%%%%%%%%%%%%%%%%%%%%%%%%%%%%%%%%
\begin{frame}
\frametitle{Exercise 1}
\label{sec-4}


\begin{itemize}
\item Find the absolute value of \texttt{3 + 4j}
\end{itemize}
\end{frame}

\begin{frame}
\frametitle{Exercise 2}
\label{sec-5}
\begin{itemize}
\item What is the datatype of number 999999999999999999? Is it not \texttt{int}?
\end{itemize}
\end{frame}
%%%%%%%%%%%%%%%%%%%%%%%%%%%%%%%%%%%%%%%%%%%%%%%%%%%%%%%%%%%%%%%%%%%%%%%%%%%%%%%%
\begin{frame}
\frametitle{Solution 2}
\label{sec-6}

\begin{itemize}
\item \texttt{Long}
\item Large integers numbers are internally stored in python as \texttt{Long}
    datatype.
\end{itemize}
\end{frame}
%%%%%%%%%%%%%%%%%%%%%%%%%%%%%%%%%%%%%%%%%%%%%%%%%%%%%%%%%%%%%%%%%%%%%%%%%%%%%%%%
\begin{frame}
\frametitle{Exercise 3}
\label{sec-7}

\begin{itemize}
\item Using python find square root of 3?
\end{itemize}
\end{frame}

\begin{frame}
\frametitle{Solution 3}
\label{sec-8}

\begin{itemize}
\item \texttt{3**0.5}
\end{itemize}
\end{frame}
%%%%%%%%%%%%%%%%%%%%%%%%%%%%%%%%%%%%%%%%%%%%%%%%%%%%%%%%%%%%%%%%%%%%%%%%%%%%%%%%
\begin{frame}
\frametitle{Exercise 4}
\label{sec-9}

\begin{itemize}
\item Is \texttt{3**1/2} and \texttt{3**0.5} same?
\end{itemize}
\end{frame}
%%%%%%%%%%%%%%%%%%%%%%%%%%%%%%%%%%%%%%%%%%%%%%%%%%%%%%%%%%%%%%%%%%%%%%%%%%%%%%%%
\begin{frame}
\frametitle{Sequence Data types}
\label{sec-10}
\begin{itemize}

\item Properties\\\pause
\label{sec-10_1}%
\begin{itemize}
\item Data in Sequence
\item Accessed using Index\pause
\end{itemize}

\item Type\\\pause
\label{sec-10_2}%
\begin{itemize}
\item \texttt{list}
\item \texttt{string}
\item \texttt{tuple}
\end{itemize}

\end{itemize} % ends low level
\end{frame}
%%%%%%%%%%%%%%%%%%%%%%%%%%%%%%%%%%%%%%%%%%%%%%%%%%%%%%%%%%%%%%%%%%%%%%%%%%%%%%%%
\begin{frame}
\frametitle{Exercise 5}
\label{sec-11}

\begin{itemize}
\item Check if \texttt{3} is an element of the list \texttt{[1,7,5,3,4]}.\\ In case it is, 
change it to \texttt{21}.
\end{itemize}

\end{frame}
%%%%%%%%%%%%%%%%%%%%%%%%%%%%%%%%%%%%%%%%%%%%%%%%%%%%%%%%%%%%%%%%%%%%%%%%%%%%%%%%
\begin{frame}
\frametitle{Exercise 6}
\label{sec-12}


\begin{itemize}
\item Convert the string ``Elizabeth is queen of england'' to ``Elizabeth is queen''
\end{itemize}

\end{frame}
%%%%%%%%%%%%%%%%%%%%%%%%%%%%%%%%%%%%%%%%%%%%%%%%%%%%%%%%%%%%%%%%%%%%%%%%%%%%%%%%
\begin{frame}
\frametitle{Summary}
\label{sec-13.1}

  In this tutorial, we have learnt to,


\begin{itemize}
\item Understand the number Datatypes -- \texttt{integer}, \texttt{float} and \texttt{complex}.
\item Know the boolean datatype and operators -- \texttt{+, *, /, **, \%} .
\item Use the sequence data types -- \texttt{list},\texttt{string} and \texttt{tuple}.
\end{itemize}
\end{frame}
%%%%%%%%%%%%%%%%%%%%%%%%%%%%%%%%%%%%%%%%%%%%%%%%%%%%%%%%%%%%%%%%%%%%%%%%%%%%%%%%
\begin{frame}
\frametitle{Summary contd..}
\label{sec-13.2}

\begin{itemize}
\item Slice sequences by using the row and column numbers.
\item Split and join a list using \texttt{split()} and \texttt{join()} function respectively.
\item Convert \texttt{string} to \texttt{tuple} and vice-versa.
\end{itemize}
\end{frame}
%%%%%%%%%%%%%%%%%%%%%%%%%%%%%%%%%%%%%%%%%%%%%%%%%%%%%%%%%%%%%%%%%%%%%%%%%%%%%%%%
\begin{frame}
\frametitle{Evaluation}
\label{sec-14}


\begin{enumerate}
\item What is the major difference between \texttt{tuples} and \texttt{lists}?\pause
\vspace{12pt}
\item Split this string on whitespaces
	\begin{itemize}
	\item \texttt{string = ``Split this string on whitespaces''}
	\end{itemize}
\end{enumerate}
\end{frame}
%%%%%%%%%%%%%%%%%%%%%%%%%%%%%%%%%%%%%%%%%%%%%%%%%%%%%%%%%%%%%%%%%%%%%%%%%%%%%%%%
\begin{frame}
\frametitle{Solutions}
\label{sec-15}


\begin{enumerate}
\item \texttt{tuples} are immutable while \texttt{lists} are not.\pause
\vspace{12pt}
\item \texttt{string.split()}
\end{enumerate}
\end{frame}
%%%%%%%%%%%%%%%%%%%%%%%%%%%%%%%%%%%%%%%%%%%%%%%%%%%%%%%%%%%%%%%%%%%%%%%%%%%%%%%%
\begin{frame}
\frametitle{FOSSEE}
{\color{blue}Free and Open-source Software for \\Science and Engineering Education} \\
\begin{itemize}
\item Goal - enabling all to use open source software tools
\item For more details, please visit {\color{blue}\url{http://fossee.in/}}
\end{itemize}
\end{frame}
%%%%%%%%%%%%%%%%%%%%%%%%%%%%%%%%%%%%%%%%%%%%%%%%%%%%%%%%%%%%%%%%%%%%%%%%%%%%%%%%
\begin{frame}
\frametitle{About the Spoken Tutorial Project}
\begin{itemize}
\item Watch the video available at {\color{blue}http://spoken-tutorial.org /What\_is\_a\_Spoken\_Tutorial}
\item It summarises the Spoken Tutorial project \pause
\item If you do not have good bandwidth, you can download and watch it
\end{itemize}
\end{frame}
%%%%%%%%%%%%%%%%%%%%%%%%%%%%%%%%%%%%%%%%%%%%%%%%%%%%%%%%%%%%%%%%%%%%%%%%%%%%%%%%
\begin{frame}
\frametitle{Spoken Tutorial Workshops}The Spoken Tutorial Project Team 
\begin{itemize}
\item Conducts workshops using spoken tutorials 
\item Gives certificates to those who pass an online test 
\item For more details, please write to \\ \hspace {0.5cm}{\color{blue}contact@spoken-tutorial.org}
\end{itemize}
\end{frame}
%%%%%%%%%%%%%%%%%%%%%%%%%%%%%%%%%%%%%%%%%%%%%%%%%%%%%%%%%%%%%%%%%%%%%%%%%%%%%%%%
\begin{frame}
\frametitle{Acknowledgements}
\begin{itemize}
\item Spoken Tutorial Project is a part of the Talk to a Teacher  project 
\item It is supported by the National Mission on Education through  ICT, MHRD, Government of India 
\item More information on this Mission is available at: \\{\color{blue}http://spoken-tutorial.org /NMEICT-Intro}
\end{itemize}
\end{frame}
%%%%%%%%%%%%%%%%%%%%%%%%%%%%%%%%%%%%%%%%%%%%%%%%%%%%%%%%%%%%%%%%%%%%%%%%%%%%%%%%
\begin{frame}

  \begin{block}{}
  \begin{center}
  \textcolor{blue}{\Large THANK YOU!} 
  \end{center}
  \end{block}
\begin{block}{}
  \begin{center}
    For more Information, visit our website\\
    \url{http://fossee.in/}
  \end{center}  
  \end{block}
\end{frame}
%%%%%%%%%%%%%%%%%%%%%%%%%%%%%%%%%%%%%%%%%%%%%%%%%%%%%%%%%%%%%%%%%%%%%%%%%%%%%%%%
\end{document}
